\chapter{复杂度、暴力法和优化来源}

复杂度是算法设计的预算。输入规模如果是 \code{n <= 100}，\complexity{O(n^3)} 可能可以接受；如果是 \code{n <= 2 \times 10^5}，通常要接近 \complexity{O(n \log n)} 或 \complexity{O(n)}；如果网格规模是 \code{m * n <= 10^5}，BFS 或 DFS 扫一遍通常合理；如果状态空间是指数级，就要考虑剪枝、记忆化或动态规划。

暴力法的作用是把问题完整展开。两数之和暴力枚举所有下标对，最长无重复子串暴力枚举所有起点和终点，岛屿数量暴力从每个未访问陆地扩展连通块。暴力方法慢，但它通常是正确性的起点。

\begin{lstlisting}[language=C++,caption={两数之和暴力解}]
std::vector<int> two_sum_bruteforce(const std::vector<int>& nums, int target)
{
    for (int left = 0; left < static_cast<int>(nums.size()); ++left) {
        for (int right = left + 1; right < static_cast<int>(nums.size()); ++right) {
            if (nums[left] + nums[right] == target) {
                return {left, right};
            }
        }
    }
    return {};
}
\end{lstlisting}

这段代码枚举的是下标对，数量约为 \code{n * (n - 1) / 2}，时间复杂度是 \complexity{O(n^2)}。它慢在每个数都要重新扫描后面的数。优化思路不是凭空出现，而是把“扫描后面所有数”改成“快速查询某个补数是否出现过”。这就是哈希表。

常见优化来源有六类。第一类是缓存重复计算，例如前缀和、记忆化搜索和动态规划。第二类是单调性，例如二分、双指针和某些滑动窗口。第三类是局部状态可维护，例如字符计数、窗口和、不同元素个数。第四类是候选淘汰，例如单调栈和单调队列。第五类是数据结构查询能力，例如哈希表、堆、平衡树、并查集。第六类是改变建模方式，例如把网格问题建成图，把区间问题排序，把字符串问题转为频次。

写完暴力法后，应该固定问几个问题：我枚举的对象是什么？有没有重复计算同一段信息？是否只需要历史信息的一小部分？是否存在有序或单调？是否可以把区间查询变成前缀差？是否可以用哈希把线性查询变成平均常数？是否可以把指数搜索合并成状态 DP？

面试时要把这些问题讲出来。一个好答案通常不是从“我用某某模板”开始，而是从暴力法、瓶颈和优化依据开始。
